% ===================================================================
%  TÁBLA Uimh. 3 — An Óige agus an Aos Óg.
%  Traduction 2026 d'apres texto/io, controlee sur texto/fr, texto/en
%  et texto/es. Consigne : voir texto/ga/10-tabla-01.tex.
%
%  Deux scenes, et l'ido subdivise beaucoup plus que le francais.
%
%  LE FAC-SIMILE OUBLIE DE METTRE « julio » EN GRAS au bloc c2-01-1.
%  La source ne bouge pas ; gravuri/korekti.json recoit sa ligne
%  irlandaise — « (Iúil nó ».
%
%  LES NOMS DES MOIS SONT LATINS ET NON CELTIQUES, sauf deux :
%  « Bealtaine » mai et « Samhain » novembre gardent le nom des fetes
%  celtiques de l'annee. Le reste — Márta, Aibreán, Meitheamh, Iúil,
%  Lúnasa — vient du latin par l'Eglise.
% ===================================================================

%%K t03-apar-1 apar
\VUsaut{3.0mm}
\VUtitre{15.0pt}{60}{TÁBLA Uimh. 3}
\VUsaut{1.6mm}
\VUtitre{11.4pt}{0}{An Óige agus an Aos Óg.}
\VUsaut{3.4mm}

%%K t03-apar-2 apar
(Tá táblaí 3 agus 4 cumtha chun stór bunúsach focal a thabhairt, i gceithre
\VUgras{radharc teaghlaigh}, ar an \VUgras{aimsir}, ar an \VUgras{aois}, ar an
\VUgras{teaghlach}, ar an \VUgras{éadach} agus ar an \VUgras{tsláinte}.)
\VUsaut{4.0mm}

%%K t03-tit-1 sub
\VUcentre{12.0pt}{0}{\textit{An Chéad Radharc.}}
\VUsaut{3.0mm}

%%K t03-tit-2 sub
\VUpk{10.2pt}{40}{Baisteadh faoin tuath.}
\VUsaut{2.4mm}

%%K t03-tit-3 sub
\VUpk{10.2pt}{40}{An tEarrach.}
\VUsaut{3.0mm}

%%K t03-c1-01-1 p
1. --- Táimid san \VUgras{Earrach}, i mí \VUgras{Márta}, \VUgras{Aibreáin} nó na
\VUgras{Bealtaine}, ar cheann de laethanta na \VUgras{seachtaine} --- \VUgras{Dé
Luain}, \VUgras{Dé Máirt} nó \VUgras{Dé Céadaoin}. Tá sé a \VUgras{deich a chlog}
ar maidin.

%%K t03-c1-02-1 p
2. --- Tá an \VUgras{aimsir} go hálainn, an \VUgras{t-aer} glan. Tá an ghrian ag
soilsiú. Tá roinnt \VUgras{scamall} beag \textsuperscript{(2)} ag ardú sa
\VUgras{spéir} ghorm \textsuperscript{(1)}.

%%K t03-c1-03-1 p
3. --- Is ar éigean a bhíonn \VUgras{duilleoga} ar na \VUgras{crainn}. Níl ar na
\VUgras{poibleoga} caola \textsuperscript{(3)} ná ar an \VUgras{bplána} ard
\textsuperscript{(17)} ach bachlóga fós.

%%K t03-c1-04-1 p
4. --- Tá na \VUgras{fáinleoga} \textsuperscript{(4)} tagtha ar ais agus tá siad ag
guairdeall le gártha áthais timpeall \VUgras{spuaice} \textsuperscript{(7)} an
\VUgras{chloigthí} \textsuperscript{(5)}, a chorónaíonn \VUgras{séipéal} an
tsráidbhaile \textsuperscript{(6)}.

%%K t03-tit-4 sub
\VUsaut{3.4mm}
\VUpk{10.2pt}{40}{An Séipéal.}
\VUsaut{3.0mm}

%%K t03-c1-05-1 p
5. --- Tá an séipéal seo an-simplí, lena \VUgras{phóirse} mór
\textsuperscript{(13)} agus lena \VUgras{chlog} mór \textsuperscript{(8)}. Tá an
\VUgras{tsnáthaid bheag} \textsuperscript{(9)} ar an uimhir X: taispeánann sí na
\VUgras{huaireanta}. Tá an \VUgras{cheann mhór} \textsuperscript{(10)} ar XII (an
meán lae); taispeánann sí na \VUgras{nóiméid}.

%%K t03-c1-06-1 p
6. --- Sa chloigtheach tá na \VUgras{cloig} \textsuperscript{(11)} ag bualadh go
meidhreach. Ar bharr an dín tá \VUgras{cros} bheag chloiche \textsuperscript{(12)}.

%%K t03-c1-07-1 p
7. --- Níl ag an séipéal an póirse mór amháin \textsuperscript{(13)}; tá doras beag
taoibh ann freisin. Trí dhoras sin tagann an \VUgras{sagart}
\textsuperscript{(15)} amach as an \VUgras{sacraistí} \textsuperscript{(14)}, ina
\VUgras{shútán}, agus téann sé i dtreo \VUgras{theach an tsagairt}
\textsuperscript{(19)}.

%%K t03-c1-08-1 p
8. --- Lena thaobh tá an \VUgras{bhean bhainne} \textsuperscript{(24)} lena
\VUgras{hasal} \textsuperscript{(23)}. Tá sí ag filleadh ón mbaile mór, mar a
dtugann sí a bainne maith do na páistí.

%%K t03-tit-5 sub
\VUsaut{4.0mm}
\VUpk{10.2pt}{40}{An tÚllord.}
\VUsaut{3.4mm}

%%K t03-c1-09-1 p
9. --- Tá an tUasal Liath (an t-asal) an-sásta leis an siúlóid mhaidine seo.
Tarraingíonn sé isteach le pléisiúr an t-aer atá cumhraithe ag an \VUgras{mbláth}
\textsuperscript{(20)} a chlúdaíonn \VUgras{crainn torthaí} \textsuperscript{(18)}
an úlloird in aice láimhe. Tá na \VUgras{péitseoga}
\textsuperscript{(18)} sin agus na \VUgras{haibreoga} sin \textsuperscript{(19)} ar
fad i mbándearg agus i mbán, agus geallann siad fómhar maith.

%%K t03-c1-10-1 p
10. --- Agus idir an dá linn tá na \VUgras{beacha} beaga \textsuperscript{(22)} ón
\VUgras{gcoirceog} \textsuperscript{(21)} ag eitilt ann chun a \VUgras{mil} mhilis
a bhailiú, ag crónán ag an obair.

%%K t03-c1-11-1 p
11. --- Ach cad a tharla? Seo chugainn páistí an tsráidbhaile ag rith agus ag
scaipeadh na \VUgras{ngéanna} scanraithe \textsuperscript{(25)}. Tá an mórshiúl seo
ag teacht amach as an séipéal tar éis \VUgras{bhaisteadh} mhac an nótaire.

%%K t03-c1-12-1 p
12. --- Dúisíonn Séamas beag ina \VUgras{chliabhán} \textsuperscript{(26)} le
bualadh meidhreach na gclog, agus iarrann a dheirfiúr Maighdlin, ar mhaith léi
freisin an mórshiúl a fheiceáil, ar a máthair í a chíoradh agus a ghléasadh ar an
tairseach.

%%K t03-c1-13-1 p
13. --- Aontaíonn an mháthair. Cuireann sí uirthi a \VUgras{ciarsúr cinn}
\textsuperscript{(28)}, a \VUgras{blús} \textsuperscript{(29)} agus a
\VUgras{naprún} \textsuperscript{(30)} agus suíonn sí ar chathaoir íseal os comhair
an dorais. Níl ar Mhaighdlin ach \VUgras{sciorta} \textsuperscript{(31)} agus
\VUgras{cabhail} \textsuperscript{(32)}.

%%K t03-c1-14-1 p
14. --- Nach sona atá sí! Déanann sí dearmad ar a bláthanna is ansa léi, a
\VUgras{cornáin} \textsuperscript{(34)}, ar a suíonn na \VUgras{féileacáin}
\textsuperscript{(35)}, a \VUgras{tiúilipí} \textsuperscript{(36)}, a
\VUgras{lilí na ngleanntán} \textsuperscript{(37)} sna \VUgras{pota}
\textsuperscript{(38)} ar an \VUgras{mballa} \textsuperscript{(33)} agus a
\VUgras{rósanna} \textsuperscript{(39)}, a dhéanann fál chomh hálainn sin. Bíonn sí
ag stánadh ar chulaithe saibhre na mban sa mhórshiúl.

%%K t03-c1-15-1 p
15. --- Is beag a bhreathnaíonn buachaillí an tsráidbhaile ar na culaithe. Ritheann
siad chun tosaigh agus brúann siad a chéile chun \VUgras{milseáin} \textsuperscript{(55)} nó
\VUgras{boinn} \textsuperscript{(52)} a fháil. Tá duine acu ag gol
\textsuperscript{(40)}.

%%K t03-c1-16-1 p
16. --- Sheas duine eile ar chrúb \VUgras{madra} \textsuperscript{(42)}, agus
teitheann seisean ag gearán agus cuireann sé \VUgras{cearc}
\textsuperscript{(43)} lena \VUgras{sicíní} \textsuperscript{(44)} ar teitheadh.

%%K t03-tit-6 sub
\VUsaut{5.0mm}
\VUpk{10.2pt}{40}{Mórshiúl an bhaiste.}
\VUsaut{4.0mm}

%%K t03-c1-17-1 p
17. --- I dtosach an mhórshiúil siúlann an \VUgras{bhuime}
\textsuperscript{(45)}. Tá \VUgras{caidhp} álainn uirthi \textsuperscript{(46)} le
ribíní fada. Iompraíonn sí an \VUgras{naíonán nuabheirthe} \textsuperscript{(47)},
ar fad i \VUgras{lása} \textsuperscript{(48)}.

%%K t03-c1-18-1 p
18. --- I ndiaidh na buime tagann an \VUgras{t-athair baiste}
\textsuperscript{(49)} agus an \VUgras{mháthair bhaiste} \textsuperscript{(53)}.
Roinneann siad milseáin agus boinn. Tá \VUgras{lámhainní}
\textsuperscript{(51)} agus \VUgras{hata ard} \textsuperscript{(50)} ar an athair
baiste. Coinníonn an mháthair bhaiste \VUgras{crobhaing bláthanna} álainn
\textsuperscript{(54)} ina láimh.

%%K t03-c1-19-1 p
19. --- Ina ndiaidh feictear \VUgras{uncail} \textsuperscript{(56)} agus
\VUgras{aintín} \textsuperscript{(58)} an linbh. Tá \VUgras{hata} galánta ar an
aintín \textsuperscript{(59)}, agus \VUgras{casóg fhada} dhubh
\textsuperscript{(57)} agus veist bhán ar an uncail. Níos faide siar tagann an
\VUgras{col ceathrar} \textsuperscript{(60)} agus an \VUgras{banchol ceathrar}
\textsuperscript{(61)}. Tugann an duine deireanach seo ar láimh \VUgras{deirfiúr}
\textsuperscript{(62)} an naíonáin, Beirtí bheag.

%%K t03-c1-20-1 p
20. --- Siúlann \VUgras{deartháir} eile Bheirtí \textsuperscript{(63)} ina ndiaidh.
Tugann sé a lámh do \VUgras{ghaol} \textsuperscript{(64)}. Ina ndiaidh siúd tagann
\VUgras{cairde} an teaghlaigh \textsuperscript{(65)}.

%%K t03-tit-7 sub
\VUsaut{4.6mm}
\VUpk{10.2pt}{40}{Lucht féachana.}
\VUsaut{3.4mm}

%%K t03-c1-21-1 p
21. --- In aice leis an mórshiúl tá \VUgras{bacach} \textsuperscript{(66)}, ag
brath ar a \VUgras{mhaide croise} \textsuperscript{(67)}. Iarrann sé déirc.

%%K t03-c1-22-1 p
22. --- Ar an taobh eile tá buachaill an-\VUgras{bhéasach}
\textsuperscript{(68)}. Ardaíonn sé a chaipín agus deir sé \VUgras{«lá maith»} leo
siúd a bhfuil aithne aige orthu.

%%K t03-c1-23-1 p
23. --- Tá muintir an tsráidbhaile tagtha amach as na tithe chun mórshiúl an
bhaiste a fheiceáil. Feicimid \VUgras{feirmeoir} \textsuperscript{(69)} le
\VUgras{caipín páipéir}, agus lena thaobh \VUgras{bean feirmeora}
\textsuperscript{(70)}. Ansin \VUgras{seanbhean} \textsuperscript{(71)}, ag brath
ar a \VUgras{bata} \textsuperscript{(72)}. Faoi dheireadh an \VUgras{muilleoir}
\textsuperscript{(73)} le \VUgras{hata leathan feilte} \textsuperscript{(74)} agus
bean feirmeora óg i \VUgras{mbróga adhmaid} \textsuperscript{(75)}, atá ag múineadh
a leanbh siúl.

%%K t03-c1-24-1 p
24. --- Is pictiúr an-bheoga é.
\VUsaut{4.0mm}

%%K t03-tit-8 sub
\VUcentre{12.0pt}{0}{\textit{An Dara Radharc.}}
\VUsaut{3.0mm}

%%K t03-tit-9 sub
\VUpk{10.2pt}{40}{Féile an phobail.}
\VUsaut{2.4mm}

%%K t03-tit-10 sub
\VUpk{10.2pt}{40}{Comórtais uisce agus rásaí.}
\VUsaut{3.0mm}

%%K t03-c2-01-1 p
1. --- Tá an \VUgras{samhradh} tagtha. Meallann grian loiscneach mhí an
\VUgras{Mheithimh} (Iúil nó \VUgras{Lúnasa}) an óige chun snámha agus chun
bádóireachta.

%%K t03-c2-02-1 p
2. --- Ar bhruach deas na habhann tá na rámhaithe ag baint díobh a gcuid éadaigh,
chun páirt a ghlacadh sna comórtais uisce agus sna rásaí.

%%K t03-c2-03-1 p
3. --- Baineann duine acu de a \VUgras{bhróga} \textsuperscript{(1)}; scaoileann sé
na hiallacha (na cnaipí). Tá beirt dá chomrádaithe réidh cheana féin. Níl orthu
anois ach \VUgras{geansaí} \textsuperscript{(2)} agus \VUgras{brístín snámha}
\textsuperscript{(3)}. Tá siad ag comhrá, ag fanacht lena gcomrádaithe. Tá siad ag
caint ar an aonach agus ar an bhféile ar an mbruach eile.

%%K t03-c2-04-1 p
4. --- Ar an talamh lena dtaobh tá \VUgras{éadaí} a gcomrádaithe: \VUgras{péire}
\VUgras{buataisí} \textsuperscript{(4)}, \VUgras{slipéir} \textsuperscript{(5)},
\VUgras{stocaí} \textsuperscript{(6)}, hataí \VUgras{feilte}
\textsuperscript{(7)} agus \VUgras{tuí} \textsuperscript{(8)} agus
\VUgras{carbhat} \textsuperscript{(9)}.

%%K t03-c2-05-1 p
5. --- Tá a \VUgras{sheaicéad} \textsuperscript{(10)} fillte go néata ag duine de
na hógánaigh. Cuireann sé ar thuáille é, ina bhfillfidh a chomharsa a
\VUgras{gculaithe} beirte gan mhoill.

%%K t03-c2-06-1 p
6. --- Tá an séú buachaill deireanach. Ar a cheann tá an \VUgras{caipín}
\textsuperscript{(11)} fós. Níor bhain sé de a \VUgras{léine}
\textsuperscript{(12)}, ná a \VUgras{bhóna} \textsuperscript{(13)}, ná a
\VUgras{chufaí} \textsuperscript{(14)}. Coinníonn sé a \VUgras{veist}
\textsuperscript{(17)} ina láimh agus tá sé fós ina \VUgras{bhríste}
\textsuperscript{(16)} agus ina \VUgras{ghealasacha} \textsuperscript{(15)}. Ní
bheidh sé in am don chéad rás eile, cinnte.

%%K t03-c2-07-1 p
7. --- In aice leis na hógánaigh tá cosán lán de \VUgras{fheochadáin}
\textsuperscript{(18)} ag dul síos go dtí an t-uisce, mar a bhfuil Séamas
críonna ag ullmhú i measc na \VUgras{ngiolcach} \textsuperscript{(19)} na
\VUgras{dtinte ealaíne} \textsuperscript{(20)}, a scaoilfear tráthnóna.

%%K t03-tit-11 sub
\VUsaut{4.4mm}
\VUpk{10.2pt}{40}{An Rás.}
\VUsaut{3.4mm}

%%K t03-c2-08-1 p
8. --- Ar an abhainn tá roinnt ban, faoi scáth a gcuid scáthanna, ag bádóireacht i
\VUgras{mbád} \textsuperscript{(21)}. Bíonn siad ag féachaint ar an rás, agus tá sé
an-taitneamhach. I ngach \VUgras{scif} \textsuperscript{(22)} bíonn ceathrar
\VUgras{rámhaithe} \textsuperscript{(24)} ag rámhaíocht lena ndícheall, agus
coinníonn an \VUgras{stiúrthóir} \textsuperscript{(26)} an \VUgras{stiúir}
\textsuperscript{(25)} agus tugann treo don árthach leochaileach. Buaileann na
\VUgras{maidí rámha} \textsuperscript{(23)} an t-uisce go rialta, agus téann na
báid éadroma ar luas meadhránach.

%%K t03-c2-09-1 p
9. --- Bíonn roinnt ógánach ag iomaíocht in aice le colún an droichid ar son duais
an \VUgras{chrainn duaise} \textsuperscript{(28)}. Bogann duine acu go cúramach ar
an gcuaille solúbtha sleamhain, chun an \VUgras{bhratach bheag}
\textsuperscript{(29)} atá ceangailte den cheann a bhaint. Tá a \VUgras{chéile
comhraic} mí-ámharach \textsuperscript{(30)} tar éis titim san uisce. Tá sé ag
snámh ar ais go dtí an bruach.

%%K t03-c2-10-1 p
10. --- Bíonn na buachaillí is sine ag \VUgras{troid ar an uisce}
\textsuperscript{(32)} in aice leis an \VUgras{gcé} \textsuperscript{(33)}. Bíonn
\VUgras{slua} mór \textsuperscript{(31)} ag féachaint ar na cluichí seo go léir, ag
brath ar \VUgras{ráille} an \VUgras{droichid} \textsuperscript{(27)} agus ag
plódú feadh an \VUgras{bhruaigh}. Casann roinnt de lucht féachana timpeall, áfach,
chun féachaint ar an gcluiche leis an \VUgras{gcuaille} \textsuperscript{(45)} nó
chun \VUgras{mórshiúl an bhardais} a mheas, atá ag dul chuig \VUgras{bronnadh na
nduaiseanna}.

%%K t03-c2-11-1 p
11. --- Ar dtús ritheann roinnt \VUgras{gasúr} \textsuperscript{(35)} os comhair na
\VUgras{gceoltóirí} \textsuperscript{(36)}; ansin tagann an \VUgras{méara}
\textsuperscript{(37)}, a chúntóirí agus \VUgras{comhairle an bhaile}. Ar deireadh
siúlann na \VUgras{comhraiceoirí dóiteáin} \textsuperscript{(38)}.

%%K t03-tit-12 sub
\VUsaut{5.0mm}
\VUpk{10.2pt}{40}{An tAonach.}
\VUsaut{3.6mm}

%%K t03-c2-12-1 p
12. --- I \VUgras{bpáirc an aonaigh} \textsuperscript{(39)} tá slua mór. Tagann na
daoine chun na \VUgras{seomraí taispeántais} éagsúla a fheiceáil: na
\VUgras{capaill adhmaid} \textsuperscript{(40)}, an \VUgras{sorcas}
\textsuperscript{(41)}, mar a bhfuil an taispeántas tosaithe cheana féin; an
\VUgras{rince poiblí} \textsuperscript{(42)}, mar a ligeann ógánaigh agus cailíní
an bhaile dóibh féin pléisiúr an \VUgras{rince}.

%%K t03-c2-13-1 p
13. --- Ar chúl feictear an \VUgras{airéine} \textsuperscript{(44)}, mar a
dtroideann \VUgras{matadóir} dána Spáinneach le \VUgras{tarbh} buile
\textsuperscript{(45)}; ansin na \VUgras{sléibhte Rúiseacha}
\textsuperscript{(49)} agus an \VUgras{raon lámhaigh} \textsuperscript{(46)}, mar a
gcleachtar \VUgras{airm thine} agus lámhach ar \VUgras{sprioc}.

%%K t03-c2-14-1 p
14. --- Lena thaobh tá \VUgras{amharclann an aonaigh} \textsuperscript{(47)}, mar a
n-imrítear \VUgras{coiméide} agus \VUgras{dráma}. An-ghar dó tá seomra an
\VUgras{fhathaigh} \textsuperscript{(48)} agus an \VUgras{abhaic}. Tá
\VUgras{airde} an chéad duine dhá mhéadar cúig déag; is ar éigean atá an dara duine
chomh mór le páiste.

%%K t03-c2-15-1 p
15. --- Faoi dheireadh, seo an \VUgras{teach ainmhithe} mór \textsuperscript{(50)}
lena \VUgras{ainmhithe fiáine}, leis an \VUgras{tíogar} \textsuperscript{(51)} agus
lena \VUgras{chrúba} cumhachtacha, leis an \VUgras{leon} \textsuperscript{(52)} agus
lena \VUgras{mhoing} tiubh, le dhá \VUgras{shioráf} \textsuperscript{(53)} agus
leis an \VUgras{eilifint} \textsuperscript{(54)}, a láimhseálann a
\VUgras{trunc} chomh haclaí sin.

%%K t03-c2-16-1 p
16. --- Tá an \VUgras{ceansaitheoir} \textsuperscript{(55)}, a thraenálann na
hainmhithe seo, ina sheasamh ag doras an tseomra.
\VUsaut{6.0mm}
\VUfilet{22mm}
